\section{Tổng quan Đề án và Tính cấp thiết}\label{sec:tong-quan}

\subsection{Bối cảnh và 3 Điểm nghẽn Cốt lõi (Pain Points)}\label{subsec:pain-points}
Trong kỷ nguyên chuyển đổi số giáo dục đại học và sau phổ thông, đặc biệt tại các lớp học có quy mô đông (từ 50 đến hơn 150 sinh viên), giảng viên phải liên tục đưa ra hàng loạt quyết định điều phối sư phạm theo thời gian thực: \textit{Nên tiếp tục giảng lý thuyết? Hay chuyển sang cho làm quiz ngắn? Đã đến lúc chia nhóm thảo luận hay cần cho cả lớp giải lao 3 phút?}

Tuy nhiên, thực tế giảng dạy hiện nay đang vấp phải ba điểm nghẽn nghiêm trọng (Hình~\ref{fig:pain-points-summary}):

\begin{enumerate}[label=\textbf{\arabic*.}]
  \item \textbf{Sự phân hóa sâu sắc của người học (Learner Divergence):} Trong cùng một tiết học, luôn tồn tại hiện tượng ``lệch pha'' năng lực: một nhóm sinh viên tiếp thu nhanh đã đạt mức nắm vững kiến thức cao ($\text{Mastery}$ cao) nhưng bắt đầu cảm thấy nhàm chán ($\text{Attention}$ suy giảm); ngược lại, nhóm sinh viên tiếp thu chậm hơn thì bị quá tải, hổng kiến thức tiên quyết và mất phương hướng. Chiến lược giảng dạy cào bằng kiểu ``một kích cỡ cho tất cả'' (\textit{One-size-fits-all}) tất yếu dẫn đến việc nhóm yếu bị bỏ lại phía sau, còn nhóm khá giỏi không được phát huy tối đa tiềm năng.
  
  \item \textbf{Điểm mù trong đánh giá trạng thái người học (Latent State Blindness):} Mức độ hiểu sâu bản chất, năng lượng nhận thức, sự mệt mỏi nhận thức (\textit{cognitive fatigue}) và mức độ gắn kết học tập (\textit{engagement}) của sinh viên đều là các \textbf{biến số trạng thái ẩn (latent states)}. Giảng viên đứng trên bục giảng hoàn toàn không thể quan sát chính xác bằng mắt thường đối với toàn bộ giảng đường rộng lớn.
  
  \item \textbf{Rủi ro khôn lường của phương pháp ``Thử và Sai'' (The Peril of Trial-and-Error):} Không một cơ sở giáo dục nào được phép biến sinh viên thật thành ``chuột bạch'' để thử nghiệm trực tiếp các thuật toán Trí tuệ Nhân tạo Tăng cường (Reinforcement Learning - RL) theo thời gian thực. Một chính sách sư phạm sai lầm từ AI khi thử nghiệm trực tiếp trên lớp học sẽ gây gián đoạn tiếp thu kiến thức, làm tổn hại tâm lý và ảnh hưởng tiêu cực đến kết quả học tập của người học.
\end{enumerate}

\begin{figure}[htbp]
  \centering
  \begin{tcolorbox}[colback=white, colframe=NavyBlue, boxrule=0.8pt, arc=3pt, width=0.98\textwidth]
    \small
    \begin{tabularx}{\textwidth}{p{3.5cm} X p{4cm}}
      \toprule
      \textbf{Điểm nghẽn (Pain Point)} & \textbf{Biểu hiện Thực tế tại Giảng đường} & \textbf{Hậu quả Sư phạm} \\
      \midrule
      \textbf{1. Phân hóa Người học} & 20\% hiểu sâu chán nản; 30\% hổng kiến thức mất phương hướng. & Áp dụng ``one-size-fits-all'' làm tăng tỷ lệ rớt môn. \\
      \addlinespace
      \textbf{2. Điểm mù Trạng thái ẩn} & Mệt mỏi nhận thức, mức gắn kết và độ hiểu bản chất không thể đo bằng mắt thường. & Giảng viên can thiệp cảm tính, thiếu căn cứ định lượng. \\
      \addlinespace
      \textbf{3. Rủi ro Thử - và - Sai} & Thử nghiệm AI trực tiếp trên sinh viên gây gián đoạn và tổn hại kết quả học tập. & Rào cản đạo đức và pháp lý ngăn cản áp dụng RL thực tế. \\
      \bottomrule
    \end{tabularx}
  \end{tcolorbox}
  \caption{Tổng hợp 3 điểm nghẽn sư phạm cốt lõi trong giảng đường truyền thống}\label{fig:pain-points-summary}
\end{figure}

\begin{highlightbox}[Giải pháp Đề xuất: Hệ thống Bản sao số Lớp học (ClassTwin)]
Để giải quyết triệt để 3 điểm nghẽn trên, đề án \textbf{ClassTwin} đề xuất một hệ thống kép đột phá:
\begin{itemize}
  \item \textbf{Classroom Digital Twin (Bản sao số Lớp học):} Mô hình hóa trạng thái ẩn của từng sinh viên và toàn bộ tập thể lớp, đồng thời đóng vai trò một môi trường giả lập (\textit{simulator}) cho phép chạy thử nghiệm hàng nghìn kịch bản phản thực (\textit{counterfactual rollout}) mà không gây bất kỳ rủi ro nào lên sinh viên thật.
  \item \textbf{Học Tăng cường Ngoại tuyến có Ràng buộc (Constrained Offline RL):} Tự động tối ưu hóa chính sách can thiệp sư phạm tối ưu từ dữ liệu lịch sử giáo dục quy mô lớn (EdNet) kết hợp ràng buộc nghiêm ngặt về giới hạn khối lượng công việc của giảng viên (\textit{workload}) và bảo đảm không làm suy giảm quyền lợi của nhóm sinh viên yếu (\textit{fairness constraint}).
  \item \textbf{Hỗ trợ Quyết định Đặt Con người ở Trung tâm (Human-in-the-Loop Decision Support):} AI đóng vai trò người đồng hành tham mưu giải thích được (\textit{explainable co-pilot}); quyền phê duyệt, tinh chỉnh hoặc từ chối quyết định can thiệp 100\% thuộc về giảng viên.
\end{itemize}
\end{highlightbox}

\subsection{Nguyên tắc Bằng chứng và Cam kết Đạo đức (Privacy-First)}\label{subsec:privacy-first}
\begin{alertbox}[Cam kết Đạo đức Dữ liệu Tuyệt đối: KHÔNG Camera Thô -- KHÔNG Ghi âm Nhận dạng]
Đề án ClassTwin kiên quyết loại bỏ hoàn toàn các phương án sử dụng camera quang học quay quét khuôn mặt (Facial Recognition/Emotion Detection) hoặc microphone thu âm định danh sinh viên. Việc phân tích cảm xúc qua camera không chỉ vi phạm nghiêm trọng quyền riêng tư học đường theo chuẩn GDPR và Luật Bảo vệ Dữ liệu Cá nhân, mà còn mang lại độ nhiễu cực lớn do hành vi đối phó của sinh viên.
Thay vào đó, ClassTwin sử dụng:
\begin{itemize}
  \item Dữ liệu hành vi học tập khách quan từ hệ thống LMS/Quiz (chuỗi đáp án, thời gian làm bài, chuỗi thao tác).
  \item Dữ liệu bối cảnh môi trường vi khí hậu vô danh tính thông qua các cảm biến IoT ($\text{CO}_2$, nhiệt độ, độ ồn decibel tổng hợp).
  \item Toàn bộ định danh sinh viên được mã hóa một chiều (\textit{salted pseudonyms}) tại Gateway lớp học trước khi truyền tải.
\end{itemize}
\end{alertbox}

\subsection{Lộ trình Triển khai 3 Cấp độ (M1 - M2 - M3)}\label{subsec:ba-cap-do}
Dự án được phân rã thành ba giai đoạn tiến hóa kỹ thuật bài bản (bảng~\ref{tab:levels-m1-m3}):
\begin{table}[htbp]
  \centering\small
  \caption{Ba cấp độ hiện thực hóa hệ thống ClassTwin}\label{tab:levels-m1-m3}
  \begin{tabularx}{\textwidth}{l p{3.2cm} X p{3.5cm}}
    \toprule
    \textbf{Cấp độ} & \textbf{Tên gọi} & \textbf{Trọng tâm Kỹ thuật} & \textbf{Tiêu chí Nghiệm thu} \\
    \midrule
    \textbf{M1} & \textbf{EdNet Baseline} & Ước lượng năng lực người học trên tập dữ liệu lịch sử chuẩn hóa EdNet (hơn 130 triệu tương tác). Dựng mô hình BKT và GRU-DKT. & Log Loss tối thiểu, ROC-AUC $\ge 0.76$, ECE $\le 0.05$ sau calibration. \\
    \addlinespace
    \textbf{M2} & \textbf{Class Simulator} & Ghép sinh viên thành lớp học ảo (40--100 người). Giả lập phản ứng phản thực đối với 5 hành động sư phạm cơ bản bằng mô hình kịch bản $\phi$. & Huấn luyện Constrained Offline RL (CQL); phân tích độ nhạy $\phi$. \\
    \addlinespace
    \textbf{M3} & \textbf{Real Pilot} & Tích hợp LMS thực tế, ứng dụng Teacher App và Gateway cảm biến IoT môi trường vào 1--3 lớp học đại học đối tác. & Độ trễ suy luận $\le 3\text{s}$, tỷ lệ giảng viên chấp thuận $\ge 65\%$. \\
    \bottomrule
  \end{tabularx}
\end{table}

\newpage

% =========================================================================
% PHẦN II: CƠ SỞ TOÁN HỌC VÀ CÔNG THỨC HÓA HỆ THỐNG
% =========================================================================
